\section{单纯同调}
\label{sec:homol-simplicial-homology}

在$\mathcal T_1$与$\mathcal T_2$中，环$0\to1\to2\to0$都是“闭合”的：沿它走一圈回到起点，没有多余的端点。区别在于，$\mathcal T_2$中这个环恰好是填充三角形的边界，而$\mathcal T_1$中没有任何二维单纯形以它为边界。要把这个区别变成可计算的量，需要能把若干条边“相加”成一个环，并对一个单纯形“取边界”。

本节固定实数系数，把同维单纯形的有限线性组合组成向量空间，称为链；边界映射是链空间之间的线性映射。闭合的链是边界映射的核，已被填掉的链是高一维边界映射的像，同调就是二者之商。这样，有限复形的孔洞计算全部归结为矩阵的核、像与秩，所用的结论（核与像是子空间、秩零度定理、商空间及其维数）可参考\BookChapterRef{chap:linear-representation}{卷二第2章“线性空间、映射与坐标”}。

\subsection{边界映射}
\label{sec:homol-boundary-map}

线段的边界是两个端点，但要让“闭合”表现为“边界为零”，端点必须带符号：从$0$走到$1$，记作“终点减起点”$1-0$；再从$1$走回$0$，记作$0-1$，两者相加为零。所以先给每个单纯形规定方向，再定义带符号的边界。

\paragraph{定向与链空间}
把$k$-单纯形$\sigma=\{v_0,\ldots,v_k\}$的顶点排成一个次序，记为$[v_0,\ldots,v_k]$，称为一个有向单纯形。同一单纯形的不同排列之间规定
\begin{equation}
 [v_{\pi(0)},\ldots,v_{\pi(k)}]=\operatorname{sgn}(\pi)\,[v_0,\ldots,v_k],
 \label{eq:homol-orientation-sign}
\end{equation}
其中$\pi$是$\{0,\ldots,k\}$的置换，$\operatorname{sgn}(\pi)$是它的符号。于是交换两个顶点改变符号，偶置换不改变它。$k\geq1$时一个单纯形恰有两种\term{定向}（orientation），彼此相差一个负号；顶点只有一种。

\begin{definition}[链空间]
\label{def:homol-chain-space}
固定$V$上的一个全序，对每个$k$-单纯形取顶点递增排列的有向单纯形。以这$n_k$个有向单纯形为基的实向量空间记为$C_k(\mathcal K)$，称为$k$维\term{链空间}（chain space），其元素称为$k$维链。$k<0$或$k$大于$\mathcal K$的维数时，规定$C_k(\mathcal K)=0$。
\end{definition}

一个$k$维链形如$\bm c=\sum_\sigma c_\sigma\sigma$，其中求和遍历$k$-单纯形，$c_\sigma\in\mathbb R$，因而$C_k(\mathcal K)\cong\mathbb R^{n_k}$。任意排列的有向单纯形都按式~\eqref{eq:homol-orientation-sign}化为基向量的$\pm1$倍，所以选哪个全序只影响基向量的符号，不影响链空间本身。系数可以是任意实数，例如$2[0,1]$表示把这条边走两遍，$-[0,1]=[1,0]$表示反向走一遍。

\paragraph{边界映射的定义}
$k$-单纯形的边界由它的$k+1$个$(k-1)$维面组成；依次去掉第$i$个顶点，并用符号$(-1)^i$交错相加。

\begin{definition}[边界映射]
\label{def:homol-boundary-map}
对$k\geq1$，在有向单纯形上令
\begin{equation}
 \partial_k[v_0,\ldots,v_k]=\sum_{i=0}^k(-1)^i[v_0,\ldots,\widehat{v_i},\ldots,v_k],
 \label{eq:homol-boundary}
\end{equation}
其中$\widehat{v_i}$表示去掉$v_i$，并线性延拓为$\partial_k:C_k(\mathcal K)\to C_{k-1}(\mathcal K)$。规定$\partial_0=0$。$\partial_k$称为\term{边界映射}（boundary map）。
\end{definition}

低维情形最直观。$\partial_1[v_0,v_1]=v_1-v_0$，即终点减起点；$\partial_2[v_0,v_1,v_2]=[v_1,v_2]-[v_0,v_2]+[v_0,v_1]=[v_0,v_1]+[v_1,v_2]+[v_2,v_0]$，是沿三角形绕行一周的三条有向边。

式~\eqref{eq:homol-boundary}对每一种排列给出一个值，必须检查它与式~\eqref{eq:homol-orientation-sign}相容，即交换顶点次序时右端也改变符号。由于任意置换都是相邻对换的复合，只需检查交换$v_j$与$v_{j+1}$。对$i\notin\{j,j+1\}$，第$i$项的面中这两个顶点仍相邻且被交换，按式~\eqref{eq:homol-orientation-sign}改变符号，而系数$(-1)^i$不变。$i=j$与$i=j+1$两项则互换：交换后去掉第$j$个位置得到的是原来去掉$v_{j+1}$的面，系数由$(-1)^{j+1}$变为$(-1)^j$，另一项同理。所以整个右端变为相反数，$\partial_k$是良好定义的线性映射。

\paragraph{边界矩阵}
在定义~\ref{def:homol-chain-space}的基下，$\partial_k$由一个$n_{k-1}\times n_k$矩阵$\bm D_k$表示，第$\sigma$列记录$\partial_k\sigma$的系数，每列恰有$k+1$个$\pm1$。

\begin{example}[两个三角形的边界矩阵]
\label{ex:homol-boundary-matrices}
在$\mathcal T_1$中按$0<1<2$定向，边的基取$a=[0,1]$、$b=[1,2]$、$c=[0,2]$。由$\partial_1a=1-0$、$\partial_1b=2-1$、$\partial_1c=2-0$，按顶点$0,1,2$排行、按$a,b,c$排列，得到
\[
 \bm D_1=\begin{bmatrix}-1&0&-1\\1&-1&0\\0&1&1\end{bmatrix}.
\]
$\mathcal T_2$还有一个二维单纯形$s=[0,1,2]$，$\partial_2s=b-c+a$，所以$\bm D_2=(1,1,-1)^\top$。直接相乘，$\bm D_1\bm D_2$的三个分量依次为$-1+0+1$、$1-1+0$、$0+1-1$，全为零。
\end{example}

$\bm D_1$的每一列在起点处为$-1$、终点处为$1$。卷三定义的关联矩阵按记录方向在起点处记$1$、终点处记$-1$，所以在相同的边方向下，$\bm D_1$恰为关联矩阵的相反数。一维边界映射只是带符号地记录了图的连接，二维以上的边界矩阵才记录填充。

\paragraph{两次取边界为零}
例~\ref{ex:homol-boundary-matrices}中$\bm D_1\bm D_2=\bm 0$不是巧合：三角形边界绕行一周回到起点，端点的贡献正负相消。这个现象对一切维数成立。

\begin{theorem}[边界的边界为零]
\label{thm:homol-boundary-squared}
对一切$k$，$\partial_{k-1}\circ\partial_k=0$。
\end{theorem}

\begin{proof}
由线性性，只需对有向单纯形$[v_0,\ldots,v_k]$验证。展开两次边界，每一项都是去掉两个顶点$v_a,v_b$（$a<b$）得到的$(k-2)$维面，并且恰好出现两次。先去掉$v_b$，它位于第$b$个位置，系数为$(-1)^b$；此时$v_a$仍位于第$a$个位置，再去掉它，系数为$(-1)^a$。先去掉$v_a$，系数为$(-1)^a$；此时$v_b$前移到第$b-1$个位置，再去掉它，系数为$(-1)^{b-1}$。两项系数分别为$(-1)^{a+b}$与$(-1)^{a+b-1}$，互为相反数，全部相消。
\end{proof}

满足$\partial_{k-1}\circ\partial_k=0$的一列向量空间与线性映射
\[
 \cdots\xrightarrow{\ \partial_{k+1}\ }C_k(\mathcal K)\xrightarrow{\ \partial_k\ }C_{k-1}(\mathcal K)\xrightarrow{\ \partial_{k-1}\ }\cdots\xrightarrow{\ \partial_1\ }C_0(\mathcal K)\xrightarrow{\ \partial_0\ }0
\]
称为\term{链复形}（chain complex）。定理~\ref{thm:homol-boundary-squared}说明，任何链的边界都是“闭合”的；下一小节把这一点写成子空间的包含关系。

\subsection{同调空间}
\label{sec:homol-homology-spaces}

边界为零的链没有多余的端点，是闭合的；某个高一维链的边界，则是已经被填掉的。定理~\ref{thm:homol-boundary-squared}说后者总是前者的一部分。“没有被填掉的闭合结构”，就是从前者中除去后者剩下的部分。

\begin{definition}[闭链、边界链与同调空间]
\label{def:homol-homology}
令
\[
 Z_k(\mathcal K)=\ker\partial_k,\qquad B_k(\mathcal K)=\operatorname{im}\partial_{k+1},
\]
其元素分别称为$k$维\term{闭链}（cycle）与$k$维\term{边界链}（boundary chain）。由定理~\ref{thm:homol-boundary-squared}，$B_k(\mathcal K)\subseteq Z_k(\mathcal K)$。商空间
\[
 H_k(\mathcal K)=Z_k(\mathcal K)/B_k(\mathcal K)
\]
称为$\mathcal K$的$k$维（实系数）单纯\term{同调空间}（homology space），其维数$\beta_k=\dim H_k(\mathcal K)$称为第$k$个\term{Betti数}（Betti number）。
\end{definition}

闭链$\bm z$所在的类$\bm z+B_k(\mathcal K)$记为$[\bm z]$，称为它的同调类。两个闭链同调，即$[\bm z]=[\bm z']$，当且仅当$\bm z-\bm z'$是某个$(k+1)$维链的边界，也就是二者相差一块被填充的部分。$[\bm z]=0$表示$\bm z$本身已被填掉。

由商空间的维数公式与秩零度定理，$\dim Z_k=n_k-\operatorname{rank}\bm D_k$，$\dim B_k=\operatorname{rank}\bm D_{k+1}$，所以
\begin{equation}
 \beta_k=n_k-\operatorname{rank}\bm D_k-\operatorname{rank}\bm D_{k+1}.
 \label{eq:homol-betti-rank}
\end{equation}
有限复形的Betti数只需数单纯形、求边界矩阵的秩。$\bm D_0$与维数以上的$\bm D_k$都按零矩阵处理。

\begin{example}[三角形边界与填实三角形的同调]
\label{ex:homol-triangle-homology}
沿用例~\ref{ex:homol-boundary-matrices}。$\bm D_1$的第三列等于前两列之和，前两列线性无关，所以$\operatorname{rank}\bm D_1=2$。解$\bm D_1\bm x=\bm0$得$x_a=x_b=-x_c$，于是
\[
 Z_1(\mathcal T_1)=Z_1(\mathcal T_2)=\operatorname{span}\{\bm z\},\qquad\bm z=a+b-c=[0,1]+[1,2]+[2,0].
\]
对$\mathcal T_1$，没有二维单纯形，$B_1=0$，由式~\eqref{eq:homol-betti-rank}，$\beta_1=3-2-0=1$，$H_1(\mathcal T_1)$由$[\bm z]$张成。对$\mathcal T_2$，$\partial_2s=\bm z$，$B_1=Z_1$，$\beta_1=3-2-1=0$；又$\bm D_2\neq\bm 0$，$\ker\partial_2=0$，$\beta_2=0$。两者都有$\beta_0=3-0-2=1$。
\end{example}

例~\ref{ex:homol-triangle-homology}把本章开头的区别变成了一个数：两个复形的一维骨架相同，闭链空间也相同，差别全在$B_1$中。框上的环是闭链而非边界，$\beta_1=1$；纸片上的同一个环是填充三角形的边界，$\beta_1=0$。

\paragraph{零维同调与连通分量}
$\beta_0$在两个例子中都是$1$，对应“只有一块”。这不是巧合。

\begin{theorem}[零维同调计数连通分量]
\label{thm:homol-h0-components}
设$\mathcal K$的一维骨架作为图有连通分量$C_1,\ldots,C_m$，从每个分量中各取一个顶点$u_j$。则$[u_1],\ldots,[u_m]$是$H_0(\mathcal K)$的一组基，特别地，$\beta_0$等于连通分量的个数。
\end{theorem}

\begin{proof}
由$\partial_0=0$，$Z_0=C_0(\mathcal K)$，$B_0$由所有边的边界$v-u$张成。若$v$与$u_j$在同一分量中，取一条路径$u_j=w_0,w_1,\ldots,w_\ell=v$，则$v-u_j=\sum_i(w_{i+1}-w_i)$是各边边界之和（反向边差一个符号），所以$[v]=[u_j]$。于是每个顶点的类都是某个$[u_j]$，$[u_1],\ldots,[u_m]$张成$H_0$。

再证线性无关。对每个分量$C_j$，定义线性函数$\varphi_j(\bm c)=\sum_{v\in C_j}c_v$，即$\bm c$在$C_j$中各顶点上系数之和。每条边的两个端点在同一分量中，所以$\varphi_j(v-u)=0$，$\varphi_j$在$B_0$上为零，因而在$H_0$上有定义。若$\sum_j\lambda_j[u_j]=0$，用$\varphi_l$作用得$\lambda_l=0$。
\end{proof}

由此也得到$\operatorname{rank}\bm D_1=n_0-m$。在例~\ref{ex:homol-triangle-homology}中，$m=1$，$\operatorname{rank}\bm D_1=3-1=2$，与直接计算一致。

\paragraph{二维的洞}
同样的计算在高维记录被围住的空腔。

\begin{example}[四面体表面的同调]
\label{ex:homol-tetrahedron-homology}
对例~\ref{ex:homol-tetrahedron-boundary}的$\mathcal S$，记去掉顶点$i$的三角形为$f_i$（顶点递增排列）。按式~\eqref{eq:homol-boundary}对四个顶点的有向单纯形形式地取边界，得到
\[
 \bm w=f_0-f_1+f_2-f_3;
\]
由定理~\ref{thm:homol-boundary-squared}的同一计算，$\partial_2\bm w=0$，所以$\bm w$是二维闭链。反过来，设$\sum_ic_if_i\in\ker\partial_2$。任意两个三角形恰好共有一条边，而每条边恰好属于两个三角形，所以这条边的系数为$c_i\pm c_j=0$，各$c_i$都由$c_0$确定，$\dim Z_2\leq1$。于是$Z_2=\operatorname{span}\{\bm w\}$，$\mathcal S$中没有三维单纯形，$\beta_2=1$，$\operatorname{rank}\bm D_2=4-1=3$。

一维骨架是四个顶点两两相邻的连通图，由定理~\ref{thm:homol-h0-components}，$\beta_0=1$，$\operatorname{rank}\bm D_1=3$。由式~\eqref{eq:homol-betti-rank}，$\beta_1=6-3-3=0$。空心四面体表面上的每个闭合环都能用表面上的三角形填掉，但整个表面围住了一个空腔。加入$\{0,1,2,3\}$后，$\bm w$成为它的边界，$\beta_2$变为$0$。
\end{example}

\paragraph{用单纯形个数核对Betti数}
Betti数各自要求秩，但它们的交错和只依赖单纯形的个数，可以用来核对手算。

\begin{theorem}[Euler示性数]
\label{thm:homol-euler-characteristic}
设$\mathcal K$的维数为$d$，则
\[
 \sum_{k=0}^d(-1)^kn_k=\sum_{k=0}^d(-1)^k\beta_k.
\]
左端称为$\mathcal K$的\term{Euler示性数}（Euler characteristic），记为$\chi(\mathcal K)$。
\end{theorem}

\begin{proof}
对$\partial_k:C_k\to C_{k-1}$用秩零度定理，$n_k=\dim Z_k+\operatorname{rank}\bm D_k$；又$\beta_k=\dim Z_k-\operatorname{rank}\bm D_{k+1}$。两式相减得$n_k-\beta_k=\operatorname{rank}\bm D_k+\operatorname{rank}\bm D_{k+1}$。乘以$(-1)^k$对$k=0,\ldots,d$求和，相邻两项的$\operatorname{rank}\bm D_{k+1}$以相反符号出现而相消，剩下$\operatorname{rank}\bm D_0$与$\operatorname{rank}\bm D_{d+1}$，二者都为零。
\end{proof}

核对：$\chi(\mathcal T_1)=3-3=0=1-1$，$\chi(\mathcal T_2)=3-3+1=1=1-0+0$，$\chi(\mathcal S)=4-6+4=2=1-0+1$。反过来，已知其余Betti数时，Euler示性数可以省去一次求秩。

\paragraph{同调记录了什么，没有记录什么}
$\beta_0$计数连通分量，$\beta_1$计数没有被填掉的独立环，$\beta_2$计数被围住的空腔。“独立”是线性代数意义上的：在$\mathcal T_1$上把$\bm z$走两遍得到$2\bm z$，它与$\bm z$属于同一条直线，不算新的洞。

同调相同不足以判定空间同胚。单点复形$\{\{0\}\}$与$\mathcal T_2$的Betti数都是$(1,0,0,\ldots)$，但单点与圆盘之间没有双射。$\mathcal T_2$与只有一条边的复形$\bigl\{\{0\},\{1\},\{0,1\}\bigr\}$也有相同的Betti数，而圆盘与线段不同胚（第~\ref{sec:homol-homotopy-equivalence}小节）。同调只记录连通与孔洞的数目，不记录维数与精细形状。实数系数还会忽略一部分只在整数系数下才能看到的信息，这里不展开。

\subsection{包含诱导的同调映射}
\label{sec:homol-induced-inclusion}

在$\mathcal T_1$上填入三角形得到$\mathcal T_2$，$\beta_1$从$1$变为$0$。我们想说的不只是“数从$1$变成$0$”，而是“原来那个环被填掉了”。为此需要把小复形的同调类送进大复形，看它变成了什么。

\paragraph{链映射}
设$\mathcal L$是$\mathcal K$的子复形。$\mathcal L$的每个单纯形也是$\mathcal K$的单纯形，按同一个全序定向，于是$C_k(\mathcal L)$的每个基向量就是$C_k(\mathcal K)$的一个基向量。这给出单射$i_k:C_k(\mathcal L)\to C_k(\mathcal K)$，把一个链原样看作大复形中的链。

$\mathcal L$对取面封闭，$\mathcal L$中单纯形的面都在$\mathcal L$中，所以在两个复形中按式~\eqref{eq:homol-boundary}取边界，得到的是同一个链：$\partial_k^{\mathcal K}\circ i_k=i_{k-1}\circ\partial_k^{\mathcal L}$。满足这种交换关系的一族线性映射有一个一般的名称。

\begin{definition}[链映射]
\label{def:homol-chain-map}
设$(C_k,\partial_k)$与$(C'_k,\partial'_k)$是两个链复形。若线性映射$f_k:C_k\to C'_k$对一切$k$满足$\partial'_k\circ f_k=f_{k-1}\circ\partial_k$，则称$(f_k)_k$为\term{链映射}（chain map）。
\end{definition}

\begin{lemma}[链映射诱导同调映射]
\label{lem:homol-chain-map-induces}
链映射$(f_k)$把闭链映为闭链、把边界链映为边界链，因而$[\bm z]\mapsto[f_k\bm z]$给出良好定义的线性映射$f_*:H_k\to H'_k$。
\end{lemma}

\begin{proof}
若$\partial_k\bm z=0$，则$\partial'_kf_k\bm z=f_{k-1}\partial_k\bm z=0$。若$\bm b=\partial_{k+1}\bm c$，则$f_k\bm b=\partial'_{k+1}f_{k+1}\bm c$。所以两个闭链相差边界链时，它们的像也相差边界链，$f_*$不依赖代表元；线性性由$f_k$的线性性得到。
\end{proof}

对子复形的包含，所得的$i_*:H_k(\mathcal L)\to H_k(\mathcal K)$称为包含\term{诱导同调映射}（induced map on homology）。由定义直接可见，恒等映射诱导恒等映射；若$\mathcal L\subseteq\mathcal K\subseteq\mathcal M$，两次包含诱导映射的复合等于直接包含的诱导映射，因为链层面上二者都是把链原样看作$\mathcal M$中的链。

\begin{example}[填入三角形使一维类变为零]
\label{ex:homol-inclusion-kills-cycle}
对$\mathcal T_1\subseteq\mathcal T_2$，$i_*$把$H_1(\mathcal T_1)$的生成元$[\bm z]$送到$[\bm z]\in H_1(\mathcal T_2)$。由于$\bm z=\partial_2s$，这个类为零，所以$i_*=0$。$H_0$上$i_*$把$[0]$送到$[0]$，是同构。
\end{example}

\paragraph{为什么要记录映射}
只比较Betti数会漏掉信息。在$\mathcal T_2$之外再放一个空心三角形，其顶点为$3,4,5$，三条边为$\{3,4\},\{4,5\},\{3,5\}$，并加一条边$\{2,3\}$把两部分连起来，得到复形$\mathcal K$。$\mathcal K$的一维骨架连通，$\beta_0=1$；$(n_0,n_1,n_2)=(6,7,1)$，$\chi(\mathcal K)=0$；$\bm D_2\neq\bm0$给出$\beta_2=0$。由定理~\ref{thm:homol-euler-characteristic}，$\beta_1=\beta_0+\beta_2-\chi=1$，$H_1(\mathcal K)$由新环$\bm z'=[3,4]+[4,5]+[5,3]$的类张成（$\bm z'$是闭链，而$B_1(\mathcal K)=\operatorname{span}\{\bm z\}$不含它）。

于是$\mathcal T_1\subseteq\mathcal K$两端的Betti数完全相同，都是$\beta_0=\beta_1=1$；但$i_*:H_1(\mathcal T_1)\to H_1(\mathcal K)$为零，原来的环被填掉了，$\mathcal K$中的环是另一个。若从$\mathcal K$中去掉二维单纯形$s$，得到复形$\mathcal K'$，则$(n_0,n_1)=(6,7)$，$\chi(\mathcal K')=-1$，$\beta_1(\mathcal K')=2$，$i_*$是单射，原来的环保留下来。比较两个复形时，要同时记录各自的同调空间以及它们之间的映射；第~\ref{chap:manifold-recovery}章的持久同调正是沿一列逐渐增大的复形追踪这些映射。
